\documentclass[a4paper,11pt,oneside]{book}
% preamble_build.tex: copy of Maximus's preamble.tex for the DE v3 build.
% The blocks marked [BUILD COPY CHANGE] are the earlier build fixes.
% The blocks marked [V3 ADDITION] are new in this version (example box,
% spacing). Everything else is unchanged from preamble.tex.

% For margins
\usepackage[a4paper, margin = 20mm]{geometry}

\usepackage[Lenny]{fncychap}
% Rejne, Glenn, Lenny

\usepackage{etoc}

% [V3 ADDITION] cancel: used to show the dx's "cancelling" in Section 11.3
\usepackage{cancel}
\renewcommand{\CancelColor}{\color{red}}

% For enumeration & math symbols
\usepackage{enumerate}
\usepackage{amsmath,amsfonts,amssymb,amsthm,graphicx,lipsum}

% For diagrams, plots and boxes
\usepackage{pgfplots}
\pgfplotsset{compat=1.18}
\usepackage{tikz}
\usepackage[most]{tcolorbox}
\tcbuselibrary{skins}
\usepackage{xcolor}
\pgfplotsset{
    standard/.style={clip = false,
    axis line style = thick,
    trig format=rad,
    axis lines = middle,
    every axis x label/.style={at={(current axis.right of origin)},anchor= west},
    every axis y label/.style={at={(current axis.above origin)},anchor=south}
    }
}

% For GC font
\usepackage{fontspec}
% [BUILD COPY CHANGE 1] The TI-84 key font is loaded only if the file exists
% (checked in ./Fonts/... and ../Fonts/...). Otherwise \ti{} falls back to a
% small boxed sans-serif key. Known letters of the TI font are mapped to their
% key names in the fallback: o = [y=], z = [mode], f = [L3].
\newcommand{\tikey}[1]{\tcbox[on line, enhanced, boxsep=0pt, left=2pt, right=2pt,
  top=1pt, bottom=1pt, boxrule=0.5pt, arc=1.5pt, colback=white,
  colframe=black!70]{\sffamily\footnotesize\strut #1}}
\ExplSyntaxOn
\NewDocumentCommand{\tifallback}{m}{
  \str_case:nnF {#1}
    { {o}{\tikey{y=}} {z}{\tikey{mode}} {f}{\tikey{L3}} }
    { \tikey{#1} }
}
\ExplSyntaxOff
\def\tifontcase{0}
\IfFileExists{Fonts/TI-84 Plus CE Key Fonts/TI84PlusCEKeys.ttf}
  {\def\tifontcase{1}}{}
\IfFileExists{../Fonts/TI-84 Plus CE Key Fonts/TI84PlusCEKeys.ttf}
  {\def\tifontcase{2}}{}
\ifcase\tifontcase\relax
  % 0: font not found, use boxed fallback keys
  \newcommand{\ti}[1]{\tifallback{#1}}
\or
  % 1: original location (as in preamble.tex)
  \newfontfamily\tifont[
    Path = {Fonts/TI-84 Plus CE Key Fonts/},
    UprightFont = {TI84PlusCEKeys.ttf}
  ]{}
  \newcommand{\ti}[1]{{\tifont #1}}
\or
  % 2: one level up (when compiling inside draft-ch1/)
  \newfontfamily\tifont[
    Path = {../Fonts/TI-84 Plus CE Key Fonts/},
    UprightFont = {TI84PlusCEKeys.ttf}
  ]{}
  \newcommand{\ti}[1]{{\tifont #1}}
\fi

% SI units
\usepackage{siunitx}

% setspace
%\usepackage{setspace}
%\onehalfspacing

% For hyperlinks within and out of document
\usepackage{hyperref}
\hypersetup{
    colorlinks = true,
    linkcolor = black,
    urlcolor = blue,
}   

% For captions for pics
\usepackage{capt-of}


% tcolorbox custom environments
% Definition ========================================================
\newtcolorbox[auto counter,number within=section]{definition}[2][drop shadow]{%
enhanced, before skip=5mm, after skip=6mm, % [V3 ADDITION] spacing
colback=blue!5!white,colframe=blue!50!black,fonttitle=\bfseries,
title= Definition~\thetcbcounter. #2,#1}

% Theorem ==========================================================
\newtcolorbox[auto counter,number within=section]{theorem}[2][drop shadow]{%
enhanced, before skip=5mm, after skip=6mm, % [V3 ADDITION] spacing
colback=black!5!white,colframe=black,fonttitle=\bfseries,
title= Theorem~\thetcbcounter. #2,#1}

% Caution box ======================================================
\newtcolorbox{caution}[1][]{enhanced,
before skip=5mm,after skip=6mm, % [V3 ADDITION] was 2mm/3mm
boxrule=1pt,left=11mm,right=2mm,top=3mm,bottom=3mm,
colback=yellow!50,
colframe=yellow!20!black,
sharp corners,rounded corners=southeast,arc is angular,arc=3mm,
underlay={%
\path[fill=tcbcolback!80!black] ([yshift=3mm]interior.south east)--++(-0.4,-0.1)--++(0.1,-0.2);
\path[draw=tcbcolframe,shorten <=-0.05mm,shorten >=-0.05mm] ([yshift=3mm]interior.south east)--++(-0.4,-0.1)--++(0.1,-0.2);
\path[fill=yellow!50!black,draw=none] (interior.south west) rectangle node[white]{\Huge\bfseries !} ([xshift=9mm]interior.north west);
},
drop fuzzy shadow,#1}

% Question box ======================================================
\newtcolorbox{question}[1][]{enhanced,
before skip=5mm,after skip=6mm, % [V3 ADDITION] was 2mm/3mm
boxrule=1pt,left=11mm,right=2mm,top=3mm,bottom=3mm,
colback=blue!20!white,
colframe=blue!20!black,
sharp corners,rounded corners=southeast,arc is angular,arc=3mm,
underlay={%
\path[fill=tcbcolback!80!black] ([yshift=3mm]interior.south east)--++(-0.4,-0.1)--++(0.1,-0.2);
\path[draw=tcbcolframe,shorten <=-0.05mm,shorten >=-0.05mm] ([yshift=3mm]interior.south east)--++(-0.4,-0.1)--++(0.1,-0.2);
\path[fill=blue!50!black,draw=none] (interior.south west) rectangle node[white]{\Huge\bfseries ?} ([xshift=9mm]interior.north west);
},
drop fuzzy shadow,#1}

% Information box ======================================================
\newtcolorbox{info}[1][]{enhanced,
before skip=5mm,after skip=6mm, % [V3 ADDITION] was 2mm/3mm
boxrule=1pt,left=11mm,right=2mm,top=3mm,bottom=3mm,
colback=cyan!20!white,
colframe=cyan!20!black,
sharp corners,rounded corners=southeast,arc is angular,arc=3mm,
underlay={%
\path[fill=tcbcolback!80!black] ([yshift=3mm]interior.south east)--++(-0.4,-0.1)--++(0.1,-0.2);
\path[draw=tcbcolframe,shorten <=-0.05mm,shorten >=-0.05mm] ([yshift=3mm]interior.south east)--++(-0.4,-0.1)--++(0.1,-0.2);
\path[fill=cyan!50!black,draw=none] (interior.south west) rectangle node[white]{\Huge\bfseries i} ([xshift=9mm]interior.north west);
},
drop fuzzy shadow,#1}

% Mistake box ======================================================
\newtcolorbox{mistake}[1][]{enhanced,
before skip=5mm,after skip=6mm, % [V3 ADDITION] was 2mm/3mm
boxrule=1pt,left=11mm,right=2mm,top=3mm,bottom=3mm,
colback=red!20!white,
colframe=red!20!black,
sharp corners,rounded corners=southeast,arc is angular,arc=3mm,
underlay={%
\path[fill=tcbcolback!80!black] ([yshift=3mm]interior.south east)--++(-0.4,-0.1)--++(0.1,-0.2);
\path[draw=tcbcolframe,shorten <=-0.05mm,shorten >=-0.05mm] ([yshift=3mm]interior.south east)--++(-0.4,-0.1)--++(0.1,-0.2);
\path[fill=red!50!black,draw=none] (interior.south west) rectangle node[white]{\Huge\bfseries $\times$} ([xshift=9mm]interior.north west);
},
drop fuzzy shadow,#1}



% [V3 ADDITION] Worked example box ==================================
% Own counter, numbered within the section (Example 11.2.1, 11.2.2, ...).
% For per-subsection numbering (11.2.1.1) change `section' to `subsection'. Blue `question' boxes are kept for thought-provoking questions
% only; worked examples always use this box. Usage:
%   \begin{example} ...question text... \end{example}
%   \solution{...}
\newtcolorbox[auto counter,number within=section]{example}[1][]{%
enhanced, breakable=false, before skip=5mm, after skip=4mm,
colback=teal!5!white, colframe=teal!55!black, boxrule=0.8pt,
left=3mm, right=3mm, top=2mm, bottom=2mm,
sharp corners, rounded corners=northwest, arc=3mm,
fonttitle=\bfseries, coltitle=white,
title=Example~\thetcbcounter, attach boxed title to top left={yshift=-2mm, xshift=3mm},
boxed title style={colback=teal!55!black, sharp corners, rounded corners=northwest, arc=2mm},
drop fuzzy shadow, #1}

% [V3 ADDITION] Spacing between sections, subsections, parts, paragraphs =
% (fncychap's chapter heading is untouched; only \section and below.)
\makeatletter
\renewcommand\section{\@startsection{section}{1}{\z@}%
  {-2ex \@plus -.5ex \@minus -.2ex}{3.2ex \@plus .3ex}%
  {\normalfont\LARGE\bfseries\color{blue!50!black}}}
\renewcommand\subsection{\@startsection{subsection}{2}{\z@}%
  {-1.5ex \@plus -.5ex \@minus -.2ex}{2.6ex \@plus .3ex}%
  {\normalfont\Large\bfseries}}
\renewcommand\subsubsection{\@startsection{subsubsection}{3}{\z@}%
  {-4ex \@plus -1ex \@minus -.5ex}{1.6ex \@plus .2ex}%
  {\normalfont\large\bfseries\color{teal!50!black}}}
\makeatother
\setlength{\parindent}{0pt}
\setlength{\parskip}{0.9ex plus 0.3ex minus 0.2ex}

% Headers and Footers
\usepackage{fancyhdr,lastpage}
% [V3 R2] grey left-margin line on the out-of-syllabus pages (item 13)
\usepackage{eso-pic}
\newcommand{\oosmargin}{\AddToShipoutPictureBG{\put(12mm,0){\color{gray!75}\rule{1.6pt}{\paperheight}}}}
\pagestyle{fancy}
\lhead{Maximus Pung Jun}
% [BUILD COPY CHANGE 2] Header bug fix: \chaptername is the word "Chapter",
% so the original header printed "Chapter 1: Chapter". \chaptermark now stores
% the bare chapter title in \leftmark, and the header uses \leftmark.
\renewcommand{\chaptermark}[1]{\markboth{#1}{}}
\rhead{Chapter \thechapter: \leftmark}
% \cfoot{Page \thepage\  of \pageref{LastPage}} % buggy, does not appear on chapter pages.

% [V3 R2] tick and cross markers for mistake boxes (G3)
\newcommand{\wrong}{\textcolor{red!70!black}{\ensuremath{\boldsymbol{\times}}}}
\newcommand{\cmark}{\textcolor{green!50!black}{\ensuremath{\checkmark}}}

\newcommand{\xmark}{\wrong}   % alias: red cross
% [V3 R2] out-of-syllabus helpers (DE notes item 13)
%   \oosmargin    : grey line in the left margin on every following page
%   \oosmarginoff : stop the margin line
%   \oosnote      : grey banner "Out of syllabus. For your interest only."
%   outofsyllabus : environment = new page + margin line + banner; ends with
%                   \clearpage and switches the margin line off again
\newcommand{\oosmarginoff}{\ClearShipoutPictureBG}
\newcommand{\oosnote}[1][Nothing in this section is examined in H2 Mathematics.]{\par\noindent\colorbox{gray!15}{\parbox{\dimexpr\textwidth-2\fboxsep}{\textbf{Out of syllabus. For your interest only.} #1}}\par}
\newenvironment{outofsyllabus}{\clearpage\oosmargin}{\clearpage\oosmarginoff}

% Notation as in Maximus's notes: italic d, display style everywhere (G1).
% Fractions inside exponents stay small, so never use these macros inside ^{}.
\newcommand{\dydx}{\dfrac{dy}{dx}}
\newcommand{\dvdx}{\dfrac{dv}{dx}}
\newcommand{\dudx}{\dfrac{du}{dx}}
\newcommand{\ddydx}{\dfrac{d^{2}y}{dx^{2}}}
\newcommand{\ddt}[1]{\dfrac{d#1}{dt}}
% Small syllabus flag (no longer used in the DE notes, kept for other chapters)
\newcommand{\syllabusflag}[1]{\par\noindent{\footnotesize\color{red!70!black}\textbf{Syllabus note.} #1}\par}


% Teacher's copy toggle ================================================
% [BUILD COPY CHANGE 3] \ifx compares the whole word (the original
% \if\conform test was always false). DE_notes_student.tex predefines
% \conform as HideSolutions.
% A student copy defines \conform as HideSolutions BEFORE loading this preamble.
\providecommand\conform{ShowSolutions}
\newcommand\ShowSolutionsKey{ShowSolutions}
% [V3 ADDITION] In the student copy the hidden solution leaves stretchable
% blank space (\vfill), so the free space on each page is shared out as
% working space below the examples instead of collecting at the bottom.
\newcommand\solution[1]{\ifx\conform\ShowSolutionsKey
  \par\noindent\textbf{Solution.} #1\par\medskip\else\par\vspace{\workspace}\fi}
% [V3 R2] minimum working space left under each example in the student copy
\newlength\workspace\setlength\workspace{3cm plus 1fill}


\newcommand{\NN}{\mathbb{N}}
\newcommand{\ZZ}{\mathbb{Z}}
\newcommand{\QQ}{\mathbb{Q}}
\newcommand{\RR}{\mathbb{R}}
\newcommand{\CC}{\mathbb{C}}


% =====================================================================
\usepackage{pifont}
% preamble_local.tex: extra packages and styles for the Vectors chapter
% ONLY. Loaded after preamble_build.tex (a verbatim copy of the shared
% book preamble from the PnC / DE build, so that it can be re-synced by
% copying it over again).
% =====================================================================
\usepackage{array, booktabs}
\usetikzlibrary{arrows.meta, calc, positioning, 3d, decorations.pathreplacing, angles, quotes}
\input{../vec_common}% vecfig.tex: shared 3D diagram style (tikz-3dplot) for Vectors I / II notes and slides.
\usepackage{tikz-3dplot}
\tdplotsetmaincoords{62}{55}
\definecolor{vA}{RGB}{31,90,180}   % first vector / line 1
\definecolor{vB}{RGB}{210,90,20}   % second vector / line 2
\definecolor{vC}{RGB}{0,128,110}   % result / third
\definecolor{vN}{RGB}{140,40,150}  % normals
\definecolor{vP}{RGB}{200,215,240} % plane 1
\definecolor{vQ}{RGB}{245,215,190} % plane 2
\definecolor{vR}{RGB}{200,235,215} % plane 3
\tikzset{
  v3/.style={tdplot_main_coords, line cap=round, line join=round, >=Stealth, scale=#1},
  v3/.default=1,
  ax/.style={->, black!50, thin},
  va/.style={->, very thick, vA},
  vb/.style={->, very thick, vB},
  vc/.style={->, very thick, vC},
  vn/.style={->, very thick, vN},
  ln/.style={very thick},
  dsh/.style={dashed, black!55},
  pp/.style={fill=vP, fill opacity=0.55, draw=vA!70!black, thin},
  pq/.style={fill=vQ, fill opacity=0.55, draw=vB!70!black, thin},
  pr/.style={fill=vR, fill opacity=0.55, draw=vC!70!black, thin},
  dot/.style={circle, fill=black, inner sep=1.2pt},
}
% \vaxes{xlen}{ylen}{zlen}
\newcommand{\vaxes}[3]{%
  \draw[ax] (0,0,0) -- (#1,0,0) node[anchor=west]{$x$};
  \draw[ax] (0,0,0) -- (0,#2,0) node[anchor=south west]{$y$};
  \draw[ax] (0,0,0) -- (0,0,#3) node[anchor=south]{$z$};
  \node[anchor=east] at (0,0,0) {\scriptsize $O$};}
% \rightmark{corner}{along1}{along2} small right-angle square
\newcommand{\rtang}[3]{\draw[thin] ($(#1)!0.25cm!(#2)$) -- ($($(#1)!0.25cm!(#2)$)+($(#1)!0.25cm!(#3)$)-(#1)$) -- ($(#1)!0.25cm!(#3)$);}
\newenvironment{vis}{\par\medskip\begin{center}}{\end{center}\medskip}
% \vang{vertex}{p1}{p2}{label}: small angle arc drawn on the page
\newcommand{\vang}[4]{%
  \coordinate (vg1) at ($(#1)!0.55cm!(#2)$); \coordinate (vg2) at ($(#1)!0.55cm!(#3)$);
  \coordinate (vgm) at ($(vg1)!0.5!(vg2)$);
  \draw[vC, thick] (vg1) .. controls ($(#1)!0.72cm!(vgm)$) .. (vg2);
  \node[vC] at ($(#1)!0.95cm!(vgm)$) {#4};}


% A little air between the parts of an example, and between lines of a list.
\setlength{\itemsep}{0.85ex}
\setlength{\parsep}{0.35ex}

% Student copy: minimum working space under each example. 2.2cm (instead of
% the book's 3cm) is meant to keep the student copy on the same page breaks
% as the teacher copy.
\setlength\workspace{2.2cm plus 1fill}

% =====================================================================
% prelim_fonts.tex: fonts of the real 9758 prelim papers (Word + MathType; pdffonts on
% ACJC 2026 P1/P2, ACJC 2025 P1, JPJC 2026, RVHS 2026, CJC/DHS/TJC/NYJC/HCI/EJC 2025):
%   text Times New Roman; cover/footer Arial; RI page footer Calibri (here Carlito);
%   maths = MathType: Times New Roman letters/digits (italic variables; upright digits,
%   function names, d, e, f), SymbolMT for the integral, sigma, pi and = + < >,
%   scripts 7 pt on 12 pt (about 58 %), tight operator spacing.
% Candidate C (teacher-approved 13 Oct 2026): the real Times New Roman and Arial
% (ttf-mscorefonts-installer) + URW Standard Symbols PS (fonts-urw-base35, a SymbolMT
% clone with the same Symbol encoding) for the integral, sigma, pi, = + < >;
% TeX Gyre Termes Math only for what MathType draws itself (radicals, stretchy brackets,
% fraction bars, arrows). House style: upright d (\mathrm{d}x), upright e and f.
% Loaded once by di_style / de_core (HW sheets) and prelim_style (Tests, solutions).
% =====================================================================
\ifdefined\prelimfontsloaded\endinput\fi
\def\prelimfontsloaded{1}
\edef\pfrestoreat{\catcode`\noexpand\@=\the\catcode`\@\relax}\makeatletter
% MathType script sizes: 7 pt superscripts on 12 pt text (58 %), 5 pt second level
\DeclareMathSizes{12}{12}{7}{5}
\DeclareMathSizes{11}{11}{6.5}{5}
\DeclareMathSizes{10.95}{10.95}{6.5}{5}
\DeclareMathSizes{10}{10}{6}{5}
\DeclareMathSizes{9}{9}{5.5}{5}
\DeclareMathSizes{14.4}{14.4}{8.5}{6}
\setmainfont{Times New Roman}
\setsansfont{Arial}
\newfontfamily\coverfont{Arial}
\newfontfamily\calibrifont{Carlito}
\usepackage{unicode-math}
\setmathfont{TeX Gyre Termes Math}
\setmathfont[range={up/{latin,Latin,num}}]{Times New Roman}
\setmathfont[range={it/{latin,Latin}}]{Times New Roman Italic}
\setmathfont[range={bfup/{latin,Latin,num}}]{Times New Roman Bold}
\setmathfont[range={"002B,"003D,"003C,"003E}]{Standard Symbols PS}
\newfontfamily\pf@sym{Standard Symbols PS}
\newfontfamily\pf@symi[FakeSlant=0.21]{Standard Symbols PS}
% MathType-like spacing (tighter than TeX around = + -)
\AtBeginDocument{\thickmuskip=2.5mu plus 1mu\medmuskip=1.5mu plus 1mu\thinmuskip=3mu
% Symbol integral and sigma (18 pt at 12 pt text, as MathType), slanted Symbol pi
 \renewcommand{\int}{\mathop{\vcenter{\hbox{\pf@sym\fontsize{\fpeval{1.5*\f@size}}{0pt}\selectfont\char"F2}}}\nolimits}%
 \renewcommand{\sum}{\mathop{\vcenter{\hbox{\pf@sym\fontsize{\fpeval{1.5*\f@size}}{0pt}\selectfont\char"E5}}}\limits}%
 \renewcommand{\pi}{\mathord{\text{\pf@symi\char"70}}}}
\pfrestoreat
% Times New Roman + MathType-style maths, as his HW docs (de-docs/common/prelim_fonts.tex)
\usepackage{needspace,longtable,multicol,enumitem,array}
\providecommand{\vv}[1]{\overrightarrow{#1}}\providecommand{\cvv}[3]{\begin{pmatrix}#1\\#2\\#3\end{pmatrix}}
\raggedbottom
\setlength{\parindent}{0pt}
\fancyhf{}\lhead{\small Vectors I: Basics Worksheet: Solutions}\rhead{\small Compiled by Maximus Pung Jun}
\cfoot{\small\thepage}
\renewcommand{\headrulewidth}{0.4pt}
\newcounter{wq}[section]
\newcommand{\Q}[1]{\par\needspace{5\baselineskip}\medskip\stepcounter{wq}\noindent\textbf{\thewq}\quad #1\par}
\newcommand{\wsec}[1]{\par\needspace{8\baselineskip}\bigskip\stepcounter{section}\noindent{\Large\scshape \thesection\quad #1}\par\medskip}
\renewcommand{\thesection}{\arabic{section}}
\begin{document}
\thispagestyle{fancy}
{\LARGE\scshape Vectors I Basics Worksheet: Solutions}\par\smallskip
{Compiled by Maximus Pung Jun}\par\smallskip
{\small\textcolor{black!60}{Basics of vectors: $\mathbf i$, $\mathbf j$, $\mathbf k$ notation, magnitude, unit vectors, parallel vectors, collinear points, the ratio theorem, the scalar product and the vector product.}}\par\medskip
\wsec{$\mathbf{i}$, $\mathbf{j}$, $\mathbf{k}$ Notation and Column Vectors}
\begin{multicols}{2}\raggedcolumns
\Q{Write $\mathbf a=3\mathbf i-2\mathbf j+5\mathbf k$ as a column vector, and write $\mathbf b=\col{4}{0}{-1}$ in terms of $\mathbf i$, $\mathbf j$ and $\mathbf k$.}
\begin{align*}\mathbf a&=\col{3}{-2}{5}\\ \mathbf b&=4\mathbf i+0\mathbf j-1\mathbf k\\ &=4\mathbf i-\mathbf k\end{align*}
\Q{The points $A$ and $B$ have coordinates $(1,2,3)$ and $(4,6,1)$. Find $\overrightarrow{AB}$ in terms of $\mathbf i$, $\mathbf j$ and $\mathbf k$.}
\begin{center}\begin{tikzpicture}[x={(0.8cm,0cm)},y={(0.4cm,0.32cm)},z={(0cm,0.8cm)},>=Stealth,line cap=round]
\draw[ax](0,0,0)--(5.8,0,0) node[right]{\scriptsize$x$};
\draw[ax](0,0,0)--(0,7,0) node[above right]{\scriptsize$y$};
\draw[ax](0,0,0)--(0,0,3.8) node[above]{\scriptsize$z$};
\node[below left] at (0,0,0){\scriptsize$O$};
\draw[dsh](1,2,3)--(1,2,0)--(1,0,0);\draw[dsh](1,2,0)--(0,2,0);
\draw[dsh](4,6,1)--(4,6,0)--(4,0,0);\draw[dsh](4,6,0)--(0,6,0);
\draw[va](0,0,0)--(1,2,3);\draw[vb](0,0,0)--(4,6,1);\draw[vc](1,2,3)--(4,6,1);
\node[dot,label=above:{\scriptsize$A(1,2,3)$}] at (1,2,3){};
\node[dot,label=right:{\scriptsize$B(4,6,1)$}] at (4,6,1){};
\end{tikzpicture}\end{center}
\begin{align*}\overrightarrow{AB}&=\overrightarrow{OB}-\overrightarrow{OA}\\ &=\col{4}{6}{1}-\col{1}{2}{3}\\ &=\col{3}{4}{-2}\\ &=3\mathbf i+4\mathbf j-2\mathbf k\end{align*}
\Q{Given $\mathbf p=2\mathbf i+\mathbf k$ and $\mathbf q=-\mathbf i+3\mathbf j$, find $\mathbf r$ such that $2\mathbf p-3\mathbf q+\mathbf r=\mathbf 0$. Give $\mathbf r$ in terms of $\mathbf i$, $\mathbf j$ and $\mathbf k$.}
\begin{align*}\mathbf p&=\col{2}{0}{1},\quad \mathbf q=\col{-1}{3}{0}\\ \mathbf r&=3\mathbf q-2\mathbf p\\ &=3\col{-1}{3}{0}-2\col{2}{0}{1}\\ &=\col{-3}{9}{0}-\col{4}{0}{2}\\ &=\col{-7}{9}{-2}\\ &=-7\mathbf i+9\mathbf j-2\mathbf k\end{align*}
\end{multicols}
\Q{For the following statements, state whether it is \textbf{true} or \textbf{false}.}
\par\smallskip{\setlength{\extrarowheight}{5pt}\begin{longtable}{|>{\centering\arraybackslash}p{0.9cm}|p{\dimexpr\linewidth-0.9cm-1.6cm-6\tabcolsep-4\arrayrulewidth\relax}|>{\centering\arraybackslash}p{1.6cm}|}
\hline &\textbf{Statement}&\textbf{T/F}\\\hline\endhead
(a)&$\mathbf i+\mathbf j+\mathbf k=\col{1}{1}{1}$\par\smallskip $\displaystyle \mathbf i+\mathbf j+\mathbf k=\col{1}{1}{1}$\rule[-6pt]{0pt}{6pt}&\textbf{T}\\\hline
(b)&The $\mathbf j$-component of $4\mathbf i-\mathbf j+2\mathbf k$ is $-1$\par\smallskip $\displaystyle 4\mathbf i-\mathbf j+2\mathbf k=\col{4}{-1}{2}$\rule[-6pt]{0pt}{6pt}&\textbf{T}\\\hline
(c)&The position vector of $P(1,-2,0)$ is $\mathbf i-2\mathbf j$\par\smallskip $\displaystyle \overrightarrow{OP}=\col{1}{-2}{0}=\mathbf i-2\mathbf j$\rule[-6pt]{0pt}{6pt}&\textbf{T}\\\hline
(d)&$3\mathbf i+2\mathbf k=\col{3}{2}{0}$\par\smallskip $\displaystyle 3\mathbf i+2\mathbf k=\col{3}{0}{2}\neq\col{3}{2}{0}$\rule[-6pt]{0pt}{6pt}&\textbf{F}\\\hline
(e)&$\overrightarrow{AB}=\overrightarrow{OA}-\overrightarrow{OB}$\par\smallskip $\displaystyle \overrightarrow{AB}=\overrightarrow{OB}-\overrightarrow{OA}$\rule[-6pt]{0pt}{6pt}&\textbf{F}\\\hline
\end{longtable}}
\clearpage\wsec{Magnitude}
\begin{multicols}{2}\raggedcolumns
\Q{Find $|2\mathbf i-3\mathbf j+6\mathbf k|$.}
\begin{align*}|2\mathbf i-3\mathbf j+6\mathbf k|&=\sqrt{2^2+(-3)^2+6^2}\\ &=\sqrt{49}\\ &=7\end{align*}
\Q{The points $A$ and $B$ have coordinates $(1,2,-1)$ and $(3,-4,2)$. Find the distance $AB$.}
\begin{align*}\overrightarrow{AB}&=\col{3}{-4}{2}-\col{1}{2}{-1}\\ &=\col{2}{-6}{3}\\ AB=\left|\overrightarrow{AB}\right|&=\sqrt{2^2+(-6)^2+3^2}\\ &=\sqrt{49}\\ &=7\end{align*}
\Q{Given $\mathbf a=\col{t}{2}{t-1}$ and $|\mathbf a|=3$, find the possible values of $t$.}
\begin{align*}|\mathbf a|^2&=t^2+2^2+(t-1)^2\\ 9&=t^2+4+t^2-2t+1\\ 0&=2t^2-2t-4\\ 0&=t^2-t-2\\ 0&=(t-2)(t+1)\\ t&=2\ \text{or}\ t=-1\end{align*}
\end{multicols}
\Q{For the following statements, state whether it is \textbf{true} or \textbf{false}.}
\par\smallskip{\setlength{\extrarowheight}{5pt}\begin{longtable}{|>{\centering\arraybackslash}p{0.9cm}|p{\dimexpr\linewidth-0.9cm-1.6cm-6\tabcolsep-4\arrayrulewidth\relax}|>{\centering\arraybackslash}p{1.6cm}|}
\hline &\textbf{Statement}&\textbf{T/F}\\\hline\endhead
(a)&$|-\mathbf a|=|\mathbf a|$ for every vector $\mathbf a$\par\smallskip $\displaystyle |-\mathbf a|=|-1|\,|\mathbf a|=|\mathbf a|$\rule[-6pt]{0pt}{6pt}&\textbf{T}\\\hline
(b)&$|3\mathbf a|=3|\mathbf a|$\par\smallskip $\displaystyle |3\mathbf a|=|3|\,|\mathbf a|=3|\mathbf a|$\rule[-6pt]{0pt}{6pt}&\textbf{T}\\\hline
(c)&$\left|\col{0}{-5}{0}\right|=-5$\par\smallskip $\displaystyle \left|\col{0}{-5}{0}\right|=\sqrt{0^2+(-5)^2+0^2}=5$\rule[-6pt]{0pt}{6pt}&\textbf{F}\\\hline
(d)&$\left|\col{1}{1}{1}\right|=1$\par\smallskip $\displaystyle \left|\col{1}{1}{1}\right|=\sqrt{1^2+1^2+1^2}=\sqrt3\neq1$\rule[-6pt]{0pt}{6pt}&\textbf{F}\\\hline
(e)&$|\mathbf i+2\mathbf j-2\mathbf k|=1+2-2=1$\par\smallskip $\displaystyle |\mathbf i+2\mathbf j-2\mathbf k|=\sqrt{1^2+2^2+(-2)^2}=3$\rule[-6pt]{0pt}{6pt}&\textbf{F}\\\hline
(f)&$|\mathbf a+\mathbf b|=|\mathbf a|+|\mathbf b|$ for all vectors $\mathbf a$, $\mathbf b$\par\smallskip $\displaystyle |\mathbf i+\mathbf j|=\sqrt2\neq|\mathbf i|+|\mathbf j|=2$\rule[-6pt]{0pt}{6pt}&\textbf{F}\\\hline
(g)&$|\mathbf a+\mathbf b|\le|\mathbf a|+|\mathbf b|$ for all vectors $\mathbf a$, $\mathbf b$\par\smallskip $\displaystyle \overrightarrow{OA}=\mathbf a,\ \overrightarrow{AB}=\mathbf b,\ \overrightarrow{OB}=\mathbf a+\mathbf b,\ OB\le OA+AB\ (\triangle OAB)$\rule[-6pt]{0pt}{6pt}&\textbf{T}\\\hline
(h)&For all $\lambda\in\mathbb R$, $|\lambda\mathbf a|=\lambda|\mathbf a|$\par\smallskip $\displaystyle \lambda=-1,\ \mathbf a=\mathbf i:\ |-\mathbf i|=1\neq-1;\ |\lambda\mathbf a|=|\lambda||\mathbf a|$\rule[-6pt]{0pt}{6pt}&\textbf{F}\\\hline
\end{longtable}}
\clearpage\wsec{Unit Vectors}
\begin{multicols}{2}\raggedcolumns
\Q{Find the unit vector in the direction of $\col{2}{-1}{2}$.}
\begin{align*}\left|\col{2}{-1}{2}\right|&=\sqrt{2^2+(-1)^2+2^2}\\ &=3\\ \text{unit vector}&=\dfrac{1}{3}\col{2}{-1}{2}\end{align*}
\Q{Find the vector of magnitude $10$ in the direction of $\col{3}{0}{-4}$.}
\begin{align*}\left|\col{3}{0}{-4}\right|&=\sqrt{3^2+0^2+(-4)^2}\\ &=5\\ \text{vector}&=10\cdot\dfrac{1}{5}\col{3}{0}{-4}\\ &=\col{6}{0}{-8}\end{align*}
\Q{The points $A$ and $B$ have coordinates $(1,0,2)$ and $(3,-6,5)$. Find the vector of magnitude $14$ in the direction \emph{opposite} to $\overrightarrow{AB}$.}
\begin{align*}\overrightarrow{AB}&=\col{3}{-6}{5}-\col{1}{0}{2}\\ &=\col{2}{-6}{3}\\ \left|\overrightarrow{AB}\right|&=\sqrt{2^2+(-6)^2+3^2}\\ &=7\\ \text{vector}&=-14\cdot\dfrac{1}{7}\col{2}{-6}{3}\\ &=\col{-4}{12}{-6}\end{align*}
\end{multicols}
\Q{For the following statements, state whether it is \textbf{true} or \textbf{false}.}
\par\smallskip{\setlength{\extrarowheight}{5pt}\begin{longtable}{|>{\centering\arraybackslash}p{0.9cm}|p{\dimexpr\linewidth-0.9cm-1.6cm-6\tabcolsep-4\arrayrulewidth\relax}|>{\centering\arraybackslash}p{1.6cm}|}
\hline &\textbf{Statement}&\textbf{T/F}\\\hline\endhead
(a)&$\col{1}{1}{1}$ is a unit vector\par\smallskip $\displaystyle \left|\col{1}{1}{1}\right|=\sqrt3\neq1$\rule[-6pt]{0pt}{6pt}&\textbf{F}\\\hline
(b)&$\dfrac13\col{1}{2}{2}$ is a unit vector\par\smallskip $\displaystyle \left|\tfrac13\col{1}{2}{2}\right|=\tfrac13\sqrt{1+4+4}=1$\rule[-6pt]{0pt}{6pt}&\textbf{T}\\\hline
(c)&$\mathbf i+\mathbf j$ is a unit vector\par\smallskip $\displaystyle |\mathbf i+\mathbf j|=\sqrt{1^2+1^2+0^2}=\sqrt2\neq1$\rule[-6pt]{0pt}{6pt}&\textbf{F}\\\hline
(d)&$|\hat{\mathbf a}|=1$ for every non-zero vector $\mathbf a$\par\smallskip $\displaystyle |\hat{\mathbf a}|=\left|\dfrac{\mathbf a}{|\mathbf a|}\right|=\dfrac{|\mathbf a|}{|\mathbf a|}=1$\rule[-6pt]{0pt}{6pt}&\textbf{T}\\\hline
(e)&$\mathbf a$ and $\hat{\mathbf a}$ are parallel\par\smallskip $\displaystyle \mathbf a=|\mathbf a|\,\hat{\mathbf a}$\rule[-6pt]{0pt}{6pt}&\textbf{T}\\\hline
\end{longtable}}
\clearpage\wsec{Parallel Vectors ($\mathbf a=\lambda\mathbf b$)}
\begin{multicols}{2}\raggedcolumns
\Q{Show that $\mathbf a=\col{2}{-1}{3}$ and $\mathbf b=\col{-6}{3}{-9}$ are parallel.}
\begin{align*}\mathbf b&=\col{-6}{3}{-9}\\ &=-3\col{2}{-1}{3}\\ &=-3\mathbf a\\ \therefore\ \mathbf a&\parallel\mathbf b\ \text{(shown)}\end{align*}
\Q{Find $p$ and $q$ such that $\col{4}{p}{-2}$ is parallel to $\col{-2}{3}{q}$.}
\begin{align*}\col{4}{p}{-2}&=\lambda\col{-2}{3}{q}\\ 4&=-2\lambda\\ \lambda&=-2\\ p&=3\lambda\\ &=-6\\ -2&=\lambda q\\ -2&=-2q\\ q&=1\end{align*}
\Q{Find all vectors of magnitude $6$ that are parallel to $\mathbf a=\col{1}{2}{-2}$.}
\begin{align*}|\mathbf a|&=\sqrt{1^2+2^2+(-2)^2}\\ &=3\\ \text{vector}&=\pm6\cdot\dfrac13\col{1}{2}{-2}\\ &=\pm\col{2}{4}{-4}\end{align*}
\end{multicols}
\Q{For the following statements, state whether it is \textbf{true} or \textbf{false}.}
\par\smallskip{\setlength{\extrarowheight}{5pt}\begin{longtable}{|>{\centering\arraybackslash}p{0.9cm}|p{\dimexpr\linewidth-0.9cm-1.6cm-6\tabcolsep-4\arrayrulewidth\relax}|>{\centering\arraybackslash}p{1.6cm}|}
\hline &\textbf{Statement}&\textbf{T/F}\\\hline\endhead
(a)&If $\mathbf a\parallel\mathbf b$, then $\mathbf a=\mathbf b$\par\smallskip $\displaystyle 2\mathbf i\parallel\mathbf i,\ 2\mathbf i\neq\mathbf i$\rule[-6pt]{0pt}{6pt}&\textbf{F}\\\hline
(b)&$\col{1}{-2}{3}$ is parallel to $\col{-2}{4}{-6}$\par\smallskip $\displaystyle \col{-2}{4}{-6}=-2\col{1}{-2}{3}$\rule[-6pt]{0pt}{6pt}&\textbf{T}\\\hline
(c)&$\col{2}{4}{6}$ is parallel to $\col{1}{2}{4}$\par\smallskip $\displaystyle \dfrac21=\dfrac42\neq\dfrac64$\rule[-6pt]{0pt}{6pt}&\textbf{F}\\\hline
(d)&$\col{3}{0}{-6}$ is parallel to $\col{-1}{0}{2}$\par\smallskip $\displaystyle \col{3}{0}{-6}=-3\col{-1}{0}{2}$\rule[-6pt]{0pt}{6pt}&\textbf{T}\\\hline
(e)&If $\mathbf a=-2\mathbf b$, then $\mathbf a$ and $\mathbf b$ are not parallel since they point in opposite directions\par\smallskip $\displaystyle \mathbf a=\lambda\mathbf b,\ \lambda=-2$\rule[-6pt]{0pt}{6pt}&\textbf{F}\\\hline
\end{longtable}}
\clearpage\wsec{Collinear Points}
\begin{multicols}{2}\raggedcolumns
\Q{Show that the points $A(1,2,3)$, $B(3,5,4)$ and $C(7,11,6)$ are collinear.}
\begin{align*}\overrightarrow{AB}&=\col{3}{5}{4}-\col{1}{2}{3}\\ &=\col{2}{3}{1}\\ \overrightarrow{BC}&=\col{7}{11}{6}-\col{3}{5}{4}\\ &=\col{4}{6}{2}\\ &=2\overrightarrow{AB}\\ \therefore\ \overrightarrow{AB}&\parallel\overrightarrow{BC},\ B\ \text{common}\\ \therefore\ A,B,C&\ \text{collinear (shown)}\end{align*}
\Q{For the points in Question 1, find the ratio $AB:BC$.}
\begin{align*}\overrightarrow{BC}&=2\overrightarrow{AB}\\ BC&=2AB\\ AB:BC&=1:2\end{align*}
\Q{The points $P(2,-1,0)$, $Q(4,1,3)$ and $R(a,7,b)$ are collinear. Find $a$ and $b$.}
\begin{align*}\overrightarrow{PR}&=\lambda\overrightarrow{PQ}\\ \col{a-2}{8}{b}&=\lambda\col{2}{2}{3}\\ 8&=2\lambda\\ \lambda&=4\\ a-2&=2(4)\\ a&=10\\ b&=3(4)\\ &=12\end{align*}
\end{multicols}
\Q{For the following statements, state whether it is \textbf{true} or \textbf{false}.}
\par\smallskip{\setlength{\extrarowheight}{5pt}\begin{longtable}{|>{\centering\arraybackslash}p{0.9cm}|p{\dimexpr\linewidth-0.9cm-1.6cm-6\tabcolsep-4\arrayrulewidth\relax}|>{\centering\arraybackslash}p{1.6cm}|}
\hline &\textbf{Statement}&\textbf{T/F}\\\hline\endhead
(a)&If $\overrightarrow{AB}\parallel\overrightarrow{CD}$, then $A$, $B$, $C$ and $D$ are collinear\par\smallskip $\displaystyle A(0,0,0),B(1,0,0),C(0,1,0),D(1,1,0):\ \overrightarrow{AB}=\overrightarrow{CD}$\rule[-6pt]{0pt}{6pt}&\textbf{F}\\\hline
(b)&If $\overrightarrow{AB}=3\overrightarrow{AC}$, then $A$, $B$ and $C$ are collinear\par\smallskip $\displaystyle \overrightarrow{AB}\parallel\overrightarrow{AC},\ A\ \text{common}$\rule[-6pt]{0pt}{6pt}&\textbf{T}\\\hline
(c)&$A(0,0,0)$, $B(1,1,1)$ and $C(2,2,3)$ are collinear\par\smallskip $\displaystyle \overrightarrow{AC}=\col{2}{2}{3}\neq\lambda\col{1}{1}{1}$\rule[-6pt]{0pt}{6pt}&\textbf{F}\\\hline
(d)&The points with position vectors $\mathbf a$, $2\mathbf a$ and $-\mathbf a$ are collinear ($\mathbf a\neq\mathbf 0$)\par\smallskip $\displaystyle 2\mathbf a-\mathbf a=\mathbf a,\ -\mathbf a-\mathbf a=-2(\mathbf a)$\rule[-6pt]{0pt}{6pt}&\textbf{T}\\\hline
\end{longtable}}
\clearpage\wsec{Ratio Theorem}
\begin{center}\begin{tikzpicture}[scale=0.85,>=Stealth]
\coordinate (O) at (0,0);\coordinate (A) at (0.6,2.4);\coordinate (B) at (5.2,1.6);
\coordinate (C) at ($(A)!0.4!(B)$);
\draw[thick,black!70](A)--(B);
\draw[va](O)--(A) node[midway,left]{\small$\mathbf a$};\draw[vb](O)--(B) node[midway,below]{\small$\mathbf b$};\draw[vc](O)--(C);
\node[dot,label=below left:$O$] at (O){};\node[dot,label=above left:$A$] at (A){};
\node[dot,label=right:$B$] at (B){};\node[dot,label=below:$C$] at (C){};
\draw[decorate,decoration={brace,amplitude=4pt,raise=3pt}](A)--(C) node[midway,above=8pt,sloped]{$\lambda$};
\draw[decorate,decoration={brace,amplitude=4pt,raise=3pt}](C)--(B) node[midway,above=8pt,sloped]{$\mu$};
\end{tikzpicture}\end{center}\noindent In the True/False statements, $C$ lies on $AB$ with $AC:CB=\lambda:\mu$, $\overrightarrow{OA}=\mathbf a$ and $\overrightarrow{OB}=\mathbf b$.
\begin{multicols}{2}\raggedcolumns
\Q{Let $\overrightarrow{OA}=\mathbf i+\mathbf j+\mathbf k$ and $\overrightarrow{OB}=\mathbf i+2\mathbf j+3\mathbf k$. Given that $C$ lies on $AB$, find the coordinates of $C$ if\par\smallskip\hangindent=2em\hangafter=1\hspace*{1em}(a)\  $AC:CB=1:2$,\par\hangindent=2em\hangafter=1\hspace*{1em}(b)\  $3AC=2CB$,\par\hangindent=2em\hangafter=1\hspace*{1em}(c)\  $2AB=5CB$,\par\hangindent=2em\hangafter=1\hspace*{1em}(d)\  Given instead that $C$ lies on $AB$ produced, and $3AC=7BC$.}
\par\medskip\noindent\begin{minipage}{\linewidth}\textbf{(a)}\begin{center}\begin{tikzpicture}[scale=0.7]
\coordinate (A) at (0,2.2);\coordinate (B) at (6,2.2);\coordinate (O) at (3,0);\coordinate (C) at (2,2.2);
\fill[black!10](O)--(A)--(B)--cycle;\draw[thick,black](O)--(A)--(B)--cycle;\draw[thick,black](O)--(C);
\foreach \p in {A,B,C,O}\fill (\p) circle (1.5pt);
\node[left] at (A){$A$};\node[right] at (B){$B$};\node[above=2pt] at (C){$C$};\node[below] at (O){$O$};
\node[above=1pt,font=\small,text=blue!70!black] at ($(A)!0.5!(C)$){$1$};\node[above=1pt,font=\small,text=blue!70!black] at ($(C)!0.5!(B)$){$2$};
\end{tikzpicture}\end{center}\begin{align*}\overrightarrow{OC}&=\dfrac{2\overrightarrow{OA}+\overrightarrow{OB}}{3}\\ &=\col{1}{\frac43}{\frac53}\\ \therefore\ C&\left(1,\tfrac43,\tfrac53\right)\end{align*}\end{minipage}\par\par\medskip\noindent\begin{minipage}{\linewidth}\textbf{(b)}\begin{center}\begin{tikzpicture}[scale=0.7]
\coordinate (A) at (0,2.2);\coordinate (B) at (6,2.2);\coordinate (O) at (3,0);\coordinate (C) at (2.4,2.2);
\fill[black!10](O)--(A)--(B)--cycle;\draw[thick,black](O)--(A)--(B)--cycle;\draw[thick,black](O)--(C);
\foreach \p in {A,B,C,O}\fill (\p) circle (1.5pt);
\node[left] at (A){$A$};\node[right] at (B){$B$};\node[above=2pt] at (C){$C$};\node[below] at (O){$O$};
\node[above=1pt,font=\small,text=blue!70!black] at ($(A)!0.5!(C)$){$2$};\node[above=1pt,font=\small,text=blue!70!black] at ($(C)!0.5!(B)$){$3$};
\end{tikzpicture}\end{center}\begin{align*}\dfrac{AC}{CB}&=\dfrac23\\ \overrightarrow{OC}&=\dfrac{2\overrightarrow{OB}+3\overrightarrow{OA}}{5}\\ &=\col{1}{\frac75}{\frac95}\\ \therefore\ C&\left(1,\tfrac75,\tfrac95\right)\end{align*}\end{minipage}\par\par\medskip\noindent\begin{minipage}{\linewidth}\textbf{(c)}\begin{center}\begin{tikzpicture}[scale=0.7]
\coordinate (A) at (0,2.2);\coordinate (B) at (6,2.2);\coordinate (O) at (3,0);\coordinate (C) at (3.6,2.2);
\fill[black!10](O)--(A)--(B)--cycle;\draw[thick,black](O)--(A)--(B)--cycle;\draw[thick,black](O)--(C);
\foreach \p in {A,B,C,O}\fill (\p) circle (1.5pt);
\node[left] at (A){$A$};\node[right] at (B){$B$};\node[above=2pt] at (C){$C$};\node[below] at (O){$O$};
\node[above=1pt,font=\small,text=blue!70!black] at ($(A)!0.5!(C)$){$3$};\node[above=1pt,font=\small,text=blue!70!black] at ($(C)!0.5!(B)$){$2$};
\end{tikzpicture}\end{center}\begin{align*}\dfrac{AB}{CB}&=\dfrac52\\ \overrightarrow{OC}&=\dfrac{3\overrightarrow{OB}+2\overrightarrow{OA}}{5}\\ &=\col{1}{\frac85}{\frac{11}5}\\ \therefore\ C&\left(1,\tfrac85,\tfrac{11}5\right)\end{align*}\end{minipage}\par\par\medskip\noindent\begin{minipage}{\linewidth}\textbf{(d)}\begin{center}\begin{tikzpicture}[scale=0.7]
\coordinate (A) at (0,2.2);\coordinate (B) at (4,2.2);\coordinate (C) at (7,2.2);\coordinate (O) at (1.6,0);
\fill[black!10](O)--(A)--(B)--cycle;\draw[thick,black](O)--(A)--(C)--cycle;\draw[thick,black](O)--(B);
\foreach \p in {A,B,C,O}\fill (\p) circle (1.5pt);
\node[left] at (A){$A$};\node[above=2pt] at (B){$B$};\node[right] at (C){$C$};\node[below] at (O){$O$};
\node[above=1pt,font=\small,text=blue!70!black] at ($(A)!0.5!(B)$){$4$};\node[above=1pt,font=\small,text=blue!70!black] at ($(B)!0.5!(C)$){$3$};
\end{tikzpicture}\end{center}\begin{align*}\dfrac{AC}{BC}&=\dfrac73\\ \overrightarrow{OB}&=\dfrac{3\overrightarrow{OA}+4\overrightarrow{OC}}{7}\\ \overrightarrow{OC}&=\dfrac{7\overrightarrow{OB}-3\overrightarrow{OA}}{4}\\ &=\col{1}{\frac{11}4}{\frac92}\\ \therefore\ C&\left(1,\tfrac{11}4,\tfrac92\right)\end{align*}\end{minipage}\par
\Q{Find the position vector of the midpoint $M$ of $AB$, where $A$ and $B$ have coordinates $(3,-1,5)$ and $(-1,5,1)$.}
\begin{align*}\overrightarrow{OM}&=\dfrac12\left[\col{3}{-1}{5}+\col{-1}{5}{1}\right]\\ &=\dfrac12\col{2}{4}{6}\\ &=\col{1}{2}{3}\end{align*}
\end{multicols}
\Q{For the following statements, state whether it is \textbf{true} or \textbf{false}.}
\par\smallskip{\setlength{\extrarowheight}{5pt}\begin{longtable}{|>{\centering\arraybackslash}p{0.9cm}|p{\dimexpr\linewidth-0.9cm-1.6cm-6\tabcolsep-4\arrayrulewidth\relax}|>{\centering\arraybackslash}p{1.6cm}|}
\hline &\textbf{Statement}&\textbf{T/F}\\\hline\endhead
(a)&$\overrightarrow{OC}=\dfrac{\mu\mathbf a+\lambda\mathbf b}{\lambda+\mu}$\par\smallskip $\displaystyle \overrightarrow{OC}=\mathbf a+\tfrac{\lambda}{\lambda+\mu}(\mathbf b-\mathbf a)=\tfrac{\mu\mathbf a+\lambda\mathbf b}{\lambda+\mu}$\rule[-6pt]{0pt}{6pt}&\textbf{T}\\\hline
(b)&$\overrightarrow{OC}=\dfrac{\lambda\mathbf a+\mu\mathbf b}{\lambda+\mu}$\par\smallskip $\displaystyle \overrightarrow{OC}=\tfrac{\mu\mathbf a+\lambda\mathbf b}{\lambda+\mu}$\rule[-6pt]{0pt}{6pt}&\textbf{F}\\\hline
(c)&If $\lambda:\mu=1:1$, then $\overrightarrow{OC}=\dfrac12(\mathbf a+\mathbf b)$\par\smallskip $\displaystyle \overrightarrow{OC}=\tfrac{1\mathbf a+1\mathbf b}{2}$\rule[-6pt]{0pt}{6pt}&\textbf{T}\\\hline
(d)&If $\lambda:\mu=1:3$, then $\overrightarrow{OC}=\dfrac{3\mathbf a+\mathbf b}{4}$\par\smallskip $\displaystyle \overrightarrow{OC}=\tfrac{3\mathbf a+1\mathbf b}{1+3}$\rule[-6pt]{0pt}{6pt}&\textbf{T}\\\hline
(e)&If $\lambda:\mu=2:1$, then $\overrightarrow{OC}=\dfrac{2\mathbf a+\mathbf b}{3}$\par\smallskip $\displaystyle \overrightarrow{OC}=\tfrac{1\mathbf a+2\mathbf b}{3}$\rule[-6pt]{0pt}{6pt}&\textbf{F}\\\hline
\end{longtable}}
\clearpage\wsec{Scalar Product}
\begingroup\setlist{itemsep=1pt,topsep=2pt,parsep=0pt}\setlength{\abovedisplayskip}{3pt}\setlength{\belowdisplayskip}{3pt}
\Q{\textbf{Properties of scalar product}\par\smallskip It is given that $\mathbf a=\cvv{1}{2}{3}$, $\mathbf b=\cvv{-3}{2}{7}$, $\mathbf c=\cvv{4}{-5}{-1}$. Find the following\par\smallskip{\setlength{\tabcolsep}{0pt}\noindent\begin{tabular}{@{\hspace{1.2em}}*{4}{p{0.225\linewidth}}}(a)\ $\mathbf a\cdot\mathbf b$ & (b)\ $\mathbf b\cdot\mathbf a$ & (c)\ $\mathbf a\cdot(\mathbf b+\mathbf c)$ & (d)\ $\mathbf a\cdot\mathbf b+\mathbf a\cdot\mathbf c$\\[4pt](e)\ $\mathbf a\cdot(2\mathbf c)$ & (f)\ $(2\mathbf a)\cdot\mathbf c$ & (g)\ $\mathbf a\cdot\mathbf a$ &\end{tabular}}\par}
\par\smallskip \textbf{Ans:} \quad (a) $22$ \quad (b) $22$ \quad (c) $13$ \quad (d) $13$ \quad (e) $-18$ \quad (f) $-18$ \quad (g) $14$
\Q{\textbf{Scalar product formula: $\mathbf a\cdot\mathbf b=|\mathbf a||\mathbf b|\cos\theta$}\par\smallskip \begin{enumerate}[label=(\alph*)]
\item Points $O$, $A$, $B$ are such that $\vv{OA}=\mathbf i+4\mathbf j-3\mathbf k$ and $\vv{OB}=5\mathbf i-\mathbf j$
\begin{enumerate}[label=(\roman*)]
\item Find the size of angle $AOB$
\item Find the size of angle $OAB$.
\end{enumerate}
\item It is given that $\mathbf c=\cvv{2}{3}{6}$, $\mathbf d=\cvv{1}{-2}{3}$ and $\mathbf e=\cvv{3}{2}{-2}$. Find the angle between the vectors $\mathbf c$ and
\begin{enumerate}[label=(\roman*)]
\item the vector $\mathbf d$
\item the vector $\mathbf e$
\item the positive $y$-axis
\end{enumerate}
\item It is given that two vectors $\mathbf a$ and $\mathbf b$ satisfy the following relationship:
\[
|\mathbf a||\mathbf b|=6.
\]
Find the value of $\mathbf a\cdot\mathbf b$ for each of the following scenarios.
\begin{enumerate}[label=(\roman*)]
\item The angle between $\mathbf a$ and $\mathbf b$ is $60^\circ$.
\item The angle between $\mathbf a$ and $\mathbf b$ is $120^\circ$.
\item Vectors $\mathbf a$ and $\mathbf b$ are perpendicular.
\end{enumerate}
\end{enumerate}}
\par\smallskip \textbf{Ans:}\par (a)(i) $87.8^\circ$ \quad (ii) $133.9^\circ$\par (b)(i) $57.7^\circ$ \quad (ii) $90^\circ$ \quad (iii) $64.6^\circ$\par (c)(i) $\mathbf a\cdot\mathbf b=3$ \quad (ii) $\mathbf a\cdot\mathbf b=-3$ \quad (iii) $\mathbf a\cdot\mathbf b=0$
\Q{\textbf{Perpendicularity of vectors}\par\smallskip \begin{enumerate}[label=(\alph*)]
\item Referred to the origin $O$, the position vectors of the points $A$ and $B$ are $\mathbf i-\mathbf j+2\mathbf k$ and $2\mathbf i+4\mathbf j+\mathbf k$ respectively. Show that $OA$ is perpendicular to $OB$.
\item The position vectors $\mathbf a$ and $\mathbf b$ are given by
\[
\mathbf a=2p\mathbf i+3p\mathbf j+6p\mathbf k\quad\text{and}\quad\mathbf b=\mathbf i-2\mathbf j+2\mathbf k
\]
where $p>0$. It is given that $|\mathbf a|=|\mathbf b|$.
\begin{enumerate}[label=(\roman*)]
\item Find the exact value of $p$.
\item Show that the vectors $(\mathbf a+\mathbf b)$ and $(\mathbf a-\mathbf b)$ are perpendicular.
\end{enumerate}
\end{enumerate}}
\par\smallskip \textbf{Ans:} (b)(i) $p=\dfrac37$
\Q{\textbf{Length of projection}\par\smallskip \begin{enumerate}[label=(\alph*)]
\item Given that $\mathbf a=2\mathbf i+\mathbf j+2\mathbf k$ and $\mathbf b=-\mathbf i+2\mathbf k-\mathbf k$, find the length of projection and the projection vector of
\begin{enumerate}[label=(\roman*)]
\item $\mathbf a$ onto $\mathbf b$,
\item $\mathbf b$ onto $\mathbf a$.
\end{enumerate}
\item With reference to the origin O, the points $A$, $B$ and $C$ have position vectors given by
\[
\mathbf a=\mathbf i+2\mathbf j-\mathbf k,\ \mathbf b=3\mathbf i+\mathbf k\ \text{and}\ \mathbf c=2\mathbf i+2\mathbf j+3\mathbf k
\]
respectively. Find the length of projection of $\vv{BA}$ onto $\vv{BC}$.
\end{enumerate}}
\par\smallskip \textbf{Ans:}\par (a)(i) $\dfrac{2}{\sqrt6}$, $\dfrac13\cvv{-1}{2}{-1}$ \qquad (ii) $\dfrac23$, $\dfrac29\cvv{2}{1}{2}$\par (b) $\dfrac23$
\endgroup
\clearpage\wsec{Vector Product}
\begingroup\setlist{itemsep=1pt,topsep=2pt,parsep=0pt}\setlength{\abovedisplayskip}{3pt}\setlength{\belowdisplayskip}{3pt}
\Q{\textbf{Properties of vector product}\par\smallskip \begin{enumerate}[label=(\alph*)]
\item Find $\cvv{3}{-1}{4}\times\cvv{5}{7}{1}$
\item Given that $\mathbf a=\cvv{1}{0}{-1}$, $\mathbf b=\cvv{1}{1}{1}$ and $\mathbf c=\cvv{-1}{1}{-1}$, find
\par\smallskip{\setlength{\tabcolsep}{0pt}\noindent\begin{tabular}{@{\hspace{1.2em}}*{4}{p{0.225\linewidth}}}(i)\ $\mathbf a\times\mathbf b$; & (ii)\ $\mathbf b\times\mathbf a$; & (iii)\ $\mathbf a\times(\mathbf b+\mathbf c)$; & (iv)\ $\mathbf a\times\mathbf b+\mathbf a\times\mathbf c$;\\[4pt](v)\ $\mathbf a\times(2\mathbf b)$; & (vi)\ $(2\mathbf a)\times\mathbf b$; & (vii)\ $\mathbf a\times\mathbf a$; &\end{tabular}}\par
\end{enumerate}}
\par\smallskip \textbf{Ans:}\par (a) $\cvv{-29}{17}{26}$\par (b)(i) $\cvv{1}{-2}{1}$ \quad (ii) $\cvv{-1}{2}{-1}$ \quad (iii) $\cvv{2}{0}{2}$ \quad (iv) $\cvv{2}{0}{2}$ \quad (v) $\cvv{2}{-4}{2}$ \quad (vi) $\cvv{2}{-4}{2}$ \quad (vii) $\mathbf 0$
\Q{\textbf{Parallel vectors}\par\smallskip \begin{enumerate}[label=(\alph*)]
\item Three non-zero vectors $\mathbf a$, $\mathbf b$ and $\mathbf c$ satisfy $\mathbf a\times3\mathbf b=2\mathbf a\times\mathbf c$. Show that $\mathbf a$ and $3\mathbf b-2\mathbf c$ are parallel.
\item Find the scalar product and vector product of the following vectors, stating whether they are parallel, perpendicular, or neither.
\begin{enumerate}[label=(\roman*)]
\item $\cvv{1}{2}{3}$ and $\cvv{4}{5}{6}$
\item $\mathbf a=\mathbf i+2\mathbf j-3\mathbf k$, $\mathbf b=3\mathbf i+\mathbf k$
\item $\cvv{1}{2}{-1}$ and $\cvv{-4}{-8}{4}$
\end{enumerate}
\end{enumerate}}
\par\smallskip \textbf{Ans:} (b) Neither; Perpendicular; Parallel
\Q{\textbf{Vector product formula}\par\smallskip Using the Vector Product, find the acute angles between the following vectors:\par\smallskip{\setlength{\tabcolsep}{0pt}\noindent\begin{tabular}{@{\hspace{1.2em}}*{3}{p{0.300\linewidth}}}(a)\ $\cvv{2}{1}{3}$ and $\cvv{4}{-1}{1}$ & (b)\ $\cvv{-1}{2}{-1}$ and $\cvv{3}{-1}{5}$ & (c)\ $\cvv{4}{3}{5}$ and $\cvv{2}{-5}{1}$\end{tabular}}\par}
\par\smallskip \textbf{Ans:} (a) $51.0^\circ$ \quad (b) $46.4^\circ$ \quad (c) $87.0^\circ$
\Q{\textbf{Perpendicular distance}\par\smallskip \begin{enumerate}[label=(\alph*)]
\item Relative to an origin $O$, two point $A$ and $B$ have position vectors given by $\mathbf a=3\mathbf i-\mathbf j+2\mathbf k$ and $\mathbf b=-2\mathbf i+3\mathbf j+2\mathbf k$. Find the following:
\begin{enumerate}[label=(\roman*)]
\item The shortest distance from $A$ to $OB$
\item The shortest distance from $B$ to $OA$
\end{enumerate}
\item Relative to the origin $O$, two points $A$ and $B$ have position vectors given by $\mathbf a=\mathbf i+\mathbf j+\mathbf k$ and $\mathbf b=\mathbf i+2\mathbf j+3\mathbf k$ respectively.
\begin{enumerate}[label=(\roman*)]
\item Find the length of projection of $\mathbf b$ onto $\mathbf a$.
\item Hence, or otherwise, find the shortest distance from $B$ to $OA$.
\end{enumerate}
\item Relative to the origin $O$, the points $A$, $B$ and $C$ have position vectors given by $\mathbf a=\mathbf i+2\mathbf j+3\mathbf k$, $\mathbf b=\mathbf i-\mathbf j+\mathbf k$ and $\mathbf c=3\mathbf i-\mathbf j+\mathbf k$. Find the shortest distance from $C$ to the line joining $A$ and $B$.
\end{enumerate}}
\par\smallskip \textbf{Ans:}\par (a)(i) $4.09$ units \quad (ii) $4.51$ units\par (b)(i) $2\sqrt3$ units \quad (ii) $\sqrt2$ units\par (c) $2$ units
\Q{\textbf{Areas of triangles \& parallelograms}\par\smallskip \begin{enumerate}[label=(\alph*)]
\item Relative to the origin $O$, the points $A$ and $B$ have position vectors $\vv{OA}=\cvv{1}{0}{3}$, $\vv{OB}=\cvv{-2}{3}{2}$. Find the area of triangle $OAB$.
\item Relative to the origin $O$, the points $A$ and $B$ have position vectors $\vv{OA}=\cvv{0}{1}{1}$, $\vv{OB}=\cvv{1}{-2}{2}$. Find the area of triangle $OAB$.
\item Relative to an origin $O$, the position vectors of the vertices $A$, $B$ and $C$ of a triangle are
\[
\mathbf i+\mathbf j-\mathbf k,\ 2\mathbf i+\mathbf k\ \text{and}\ -\mathbf i+2\mathbf j-\mathbf k
\]
respectively. By considering the displacement vectors $\vv{AB}$ and $\vv{AC}$, find the exact area of triangle $ABC$.
\end{enumerate}}
\par\smallskip \textbf{Ans:} \quad (a) $\dfrac{\sqrt{154}}{2}$ units$^2$ \quad (b) $\dfrac{3\sqrt2}{2}$ units$^2$ \quad (c) $\dfrac{\sqrt{21}}{2}$ units$^2$
\endgroup
\clearpage\wsec{Practice Problems}
\Q{\textcolor{black!60}{\small[RI 2026 Prelim P1 Q4]}\par\smallskip With reference to the origin $O$, three distinct points $A$, $B$ and $C$ are such that $\overrightarrow{OA}=\mathbf a$, $\overrightarrow{OB}=\mathbf b$ and $\overrightarrow{OC}=\mathbf c$. The mid-point of $AC$ is $D$.\begin{enumerate}[label=\textbf{(\alph*)},leftmargin=2.2em,itemsep=3pt]\item Express $\overrightarrow{OD}$ in terms of $\mathbf a$ and $\mathbf c$.\hfill[1]\item Given that $|\mathbf a|=|\mathbf c|$, $|\mathbf b|\neq0$ and $BD$ is perpendicular to $AC$, show that $OB$ is also perpendicular to $AC$ by using a suitable scalar product.\hfill[3]\item Given further that point $D$ lies on line segment $OB$, $|\mathbf a|=2$, $|\mathbf b|=6$ and $\angle AOC=60^\circ$, find the exact area of quadrilateral $OABC$.\hfill[2]\end{enumerate}}
\par\smallskip\noindent\textbf{Solution}\par \par\smallskip\noindent\textbf{(a)}\par\nopagebreak{}\begin{align*}\overrightarrow{OD}&=\frac{\mathbf a+\mathbf c}{2}\end{align*}
\par\smallskip\noindent\textbf{(b)}\par\nopagebreak{}\begin{align*}\overrightarrow{BD}&=\frac{\mathbf a+\mathbf c}{2}-\mathbf b\\ &\text{Since $BD$ is perpendicular to $AC$,}\\ \overrightarrow{BD}\cdot\overrightarrow{AC}&=0\\ \Rightarrow\left(\frac{\mathbf a+\mathbf c}{2}-\mathbf b\right)\cdot(\mathbf c-\mathbf a)&=0\\ \Rightarrow\tfrac12\mathbf a\cdot\mathbf c-\tfrac12\mathbf a\cdot\mathbf a+\tfrac12\mathbf c\cdot\mathbf c-\tfrac12\mathbf c\cdot\mathbf a-\mathbf b\cdot\mathbf c+\mathbf b\cdot\mathbf a&=0\\ \Rightarrow\tfrac12\mathbf a\cdot\mathbf c-\tfrac12|\mathbf a|^2+\tfrac12|\mathbf c|^2-\tfrac12\mathbf a\cdot\mathbf c-\mathbf b\cdot\mathbf c+\mathbf b\cdot\mathbf a&=0\\ \Rightarrow-\mathbf b\cdot\mathbf c+\mathbf b\cdot\mathbf a&=0\\ \Rightarrow\mathbf b\cdot(\mathbf a-\mathbf c)&=0\end{align*}Since $\mathbf b\neq\mathbf 0$ and $\mathbf a\neq\mathbf c$, $OB$ is perpendicular to $AC$. (shown)
\par\smallskip\noindent\textbf{(c)}\par\nopagebreak{}Given that point $D$ lies on line segment $OB$, then $OB$ (or $DB$) is the perpendicular bisector of $AC$.\begin{center}\begin{tikzpicture}[scale=0.55]\coordinate (O) at (0,0);\coordinate (B) at (0,5);\coordinate (A) at (-1.2,2.08);\coordinate (C) at (1.2,2.08);\coordinate (D) at (0,2.08);\draw (O)--(A)--(B)--(C)--cycle;\draw (O)--(B);\draw (A)--(C);\node[below] at (O){\small$O$};\node[above] at (B){\small$B$};\node[left] at (A){\small$A$};\node[right] at (C){\small$C$};\node[below right] at (D){\small$D$};\end{tikzpicture}\end{center}Area of $\Delta OAB=$ Area of $\Delta OCB$\begin{align*}\text{Area of }OABC&=2\times\text{Area of }\Delta OAB\\ &=2\left(\tfrac12|\mathbf a\times\mathbf b|\right)\\ &=2\times\tfrac12|\mathbf a||\mathbf b|\sin\left(\tfrac{60}{2}\right)^\circ\\ &=2\times6\times\tfrac12=6\text{ units}^2\end{align*}
\par\medskip\hrule\medskip
\Q{\textcolor{black!60}{\small[YIJC 2026 Prelim P2 Q2]}\par\smallskip \begin{enumerate}[label=\textbf{(\alph*)},leftmargin=2.2em,itemsep=3pt]\item The non-zero vectors $\mathbf a$, $\mathbf b$ and $\mathbf c$ are such that $\mathbf a\times\mathbf c$ is parallel to $\mathbf a\times\mathbf b$. Show that $\mathbf c=m\mathbf a+n\mathbf b$, where $m$ and $n$ are constants.\hfill[3]\item Given that $\mathbf a$ is a non-zero vector and $\mathbf b$ is a unit vector such that $|\mathbf a\times\mathbf b|=2$ and $|\mathbf a\cdot\mathbf b|=2.5$, show that the acute angle between the vectors $\mathbf a$ and $\mathbf b$ is $\tan^{-1}\left(\dfrac45\right)$ and find the exact value of $|\mathbf a|$.\hfill[4]\end{enumerate}}
\par\smallskip\noindent\textbf{Solution}\par \par\smallskip\noindent\textbf{(a)}\par\nopagebreak{}\begin{align*}\mathbf a\times\mathbf c&=n(\mathbf a\times\mathbf b)\ \text{ for some (non-zero) constant }n\\ \mathbf a\times\mathbf c&=\mathbf a\times n\mathbf b\\ \mathbf a\times\mathbf c-\mathbf a\times n\mathbf b&=\mathbf 0\\ \mathbf a\times(\mathbf c-n\mathbf b)&=\mathbf 0\end{align*}Given non-zero $\mathbf a$,\par $\mathbf c-n\mathbf b=\mathbf 0$ or $\mathbf c-n\mathbf b\parallel\mathbf a$\par $\mathbf c-n\mathbf b=\mathbf 0$ or $\mathbf c-n\mathbf b=m\mathbf a$ for some constant $m\neq0$\par Hence $\mathbf c=m\mathbf a+n\mathbf b$, $m$ and $n$ are constants.
\par\smallskip\noindent\textbf{(b)}\par\nopagebreak{}\begin{center}\begin{tikzpicture}[scale=0.8,>=Stealth]\draw[thick,->](0,0)--(4,2.4) node[midway,above left]{\small$\mathbf a$};\draw[thick,->](0,0)--(0.8,0) node[below]{\small$\mathbf b$};\draw[dashed](0,0)--(4,0)--(4,2.4);\draw (3.8,0)--(3.8,0.2)--(4,0.2);\draw (0.7,0) arc(0:31:0.7);\node at (1.05,0.25){\scriptsize$\theta$};\draw[decorate,decoration={brace,amplitude=4pt}](4.15,2.4)--(4.15,0) node[midway,right=4pt]{\small$|\mathbf a\times\mathbf b|$};\draw[decorate,decoration={brace,amplitude=4pt,mirror}](0,-0.45)--(4,-0.45) node[midway,below=4pt]{\small$|\mathbf a\cdot\mathbf b|$};\end{tikzpicture}\end{center}Let $\theta$ be the acute angle between vectors $\mathbf a$ and $\mathbf b$\begin{align*}\tan\theta&=\frac{2}{2.5}\\ \theta&=\tan^{-1}\left(\frac45\right)\ \text{(shown)}\\ |\mathbf a|&=\sqrt{2^2+2.5^2}=\sqrt{\frac{41}{4}}=\frac{\sqrt{41}}{2}\end{align*}
\par\medskip\hrule\medskip
\Q{\textcolor{black!60}{\small[RI 2025 Prelim P1 Q3]}\par\smallskip Relative to the origin $O$, the points $A$ and $C$ have position vectors $\mathbf a$ and $2\mathbf b$ such that \[\mathbf a=5p\mathbf i-2p\mathbf j+4p\mathbf k\ \text{ and }\ \mathbf b=\mathbf i-2\mathbf j+2\mathbf k,\] where $p$ is a positive constant. It is given that $|\mathbf a|=2|\mathbf b|$.\begin{enumerate}[label=\textbf{(\alph*)},leftmargin=2.2em,itemsep=3pt]\item Find the exact value of $p$.\hfill[2]\item Evaluate $(\mathbf a+2\mathbf b)\cdot(\mathbf a-2\mathbf b)$.\hfill[1]\item Evaluate $\mathbf a\times\mathbf b$ and hence find the area of triangle $OAC$.\hfill[3]\item Use a geometrical reason to explain why $|\mathbf a\times\mathbf b|=\dfrac14\left(|\mathbf a+2\mathbf b|\right)\left(|\mathbf a-2\mathbf b|\right)$.\hfill[2]\end{enumerate}}
\par\smallskip\noindent\textbf{Solution}\par \par\smallskip\noindent\textbf{(a)}\par\nopagebreak{}Since $|\mathbf a|=2|\mathbf b|$,\begin{align*}\sqrt{(5p)^2+(-2p)^2+(4p)^2}&=2\sqrt{1^2+(-2)^2+2^2}\\ 3\sqrt5p&=6\qquad\text{(given }p>0)\\ p&=\frac{2}{\sqrt5}=\frac{2\sqrt5}{5}\end{align*}
\par\smallskip\noindent\textbf{(b)}\par\nopagebreak{}\begin{align*}&(\mathbf a+2\mathbf b)\cdot(\mathbf a-2\mathbf b)\\ &=\mathbf a.\mathbf a-2\mathbf a.\mathbf b+2\mathbf b.\mathbf a-4\mathbf b.\mathbf b\\ &=|\mathbf a|^2-4|\mathbf b|^2\quad\text{since }|\mathbf a|=2|\mathbf b|\\ &=0\end{align*}
\par\smallskip\noindent\textbf{(c)}\par\nopagebreak{}\begin{align*}\mathbf a\times\mathbf b=\begin{pmatrix}5p\\-2p\\4p\end{pmatrix}\times\begin{pmatrix}1\\-2\\2\end{pmatrix}=p\begin{pmatrix}5\\-2\\4\end{pmatrix}\times\begin{pmatrix}1\\-2\\2\end{pmatrix}=p\begin{pmatrix}4\\-6\\-8\end{pmatrix}=\frac{4\sqrt5}{5}\begin{pmatrix}2\\-3\\-4\end{pmatrix}\end{align*}Area of triangle $OAC$\begin{align*}&=\tfrac12|\mathbf a\times2\mathbf b|\\ &=|\mathbf a\times\mathbf b|\\ &=\frac{4\sqrt5}{5}\sqrt{2^2+(-3)^2+(-4)^2}=\frac{4\sqrt{145}}{5}\end{align*}
\par\smallskip\noindent\textbf{(d)}\par\nopagebreak{}Observe that $\overrightarrow{OD}=\mathbf a+2\mathbf b$ and $\overrightarrow{CA}=\mathbf a-2\mathbf b$.\par From \textbf{(b)}, $OD$ is perpendicular to $CA$.\par In addition, $CA$ and $OD$, which are the diagonals of the parallelogram $OADC$, bisect each other.\begin{center}\begin{tikzpicture}[scale=0.7,>=Stealth]\coordinate (O) at (0,0);\coordinate (A) at (3.6,0.2);\coordinate (C) at (1.6,3.1);\coordinate (D) at (5.2,3.3);\draw[green!50!black,thick] (O)--(A)--(D)--(C)--cycle;\draw (O)--(D);\draw (C)--(A);\draw[green!50!black,->,thick](O)--($(O)!0.55!(A)$);\draw[green!50!black,->,thick](O)--($(O)!0.55!(C)$);\node[below left] at (O){\small$O$};\node[below] at (A){\small$A$};\node[above left] at (C){\small$C$};\node[above right] at (D){\small$D$};\node[below] at ($(O)!0.5!(A)$){\small$\mathbf a$};\node[left] at ($(O)!0.5!(C)$){\small$2\mathbf b$};\end{tikzpicture}\end{center}From \textbf{(c)}, area of triangle $OAC=|\mathbf a\times\mathbf b|$.\par But area of triangle $OAC$ is also\begin{align*}&=\tfrac12(CA)\left(\tfrac12OD\right)\\ &=\tfrac14\left(|\mathbf a-2\mathbf b|\right)\left(|\mathbf a+2\mathbf b|\right)\end{align*}Hence $|\mathbf a\times\mathbf b|=\frac14\left(|\mathbf a+2\mathbf b|\right)\left(|\mathbf a-2\mathbf b|\right)$.
\par\medskip\hrule\medskip
\Q{\textcolor{black!60}{\small[ACJC 2025 Promo Q7]}\par\smallskip Referred to the origin $O$, points $A$, $B$ and $C$ have position vectors $\mathbf a$, $\mathbf b$ and $\mathbf c$ respectively. It is given that $\mathbf a+\mathbf b-\mathbf c$ and $\mathbf a-\mathbf b+\mathbf c$ are parallel.\begin{enumerate}[label=\textbf{(\alph*)},leftmargin=2.2em,itemsep=3pt]\item Using vector product, show that $\overrightarrow{OA}$ and $\overrightarrow{BC}$ are parallel.\hfill[4]\end{enumerate}It is now given that $\mathbf a$ and $\mathbf b$ are unit vectors, $\overrightarrow{BC}=2\overrightarrow{OA}$ and angle $OBC=60^\circ$.\begin{enumerate}[label=\textbf{(\alph*)},leftmargin=2.2em,itemsep=3pt,start=2]\item Using scalar product, show that $\overrightarrow{OB}$ and $\overrightarrow{OC}$ are perpendicular.\hfill[3]\item The point $H$ has position vector $\mathbf h$. Given that $\mathbf h=\mathbf b\times\mathbf c$ and $|\mathbf c|=\sqrt3$, find the exact volume of the pyramid $HOBC$.\hfill[3]\par[The volume of a pyramid is $\frac13\times$ base area $\times$ height.]\end{enumerate}}
\par\smallskip\noindent\textbf{Solution}\par \par\smallskip\noindent\textbf{(a)}\par\nopagebreak{}Since $\mathbf a+\mathbf b-\mathbf c$ and $\mathbf a-\mathbf b+\mathbf c$ are parallel,\begin{align*}(\mathbf a+\mathbf b-\mathbf c)\times(\mathbf a-\mathbf b+\mathbf c)&=\mathbf 0\\ (\mathbf a-(\mathbf c-\mathbf b))\times(\mathbf a+(\mathbf c-\mathbf b))&=\mathbf 0\\ \left(\mathbf a-\overrightarrow{BC}\right)\times\left(\mathbf a+\overrightarrow{BC}\right)&=\mathbf 0\\ (\mathbf a\times\mathbf a)+\left(\mathbf a\times\overrightarrow{BC}\right)+\left(-\overrightarrow{BC}\times\mathbf a\right)+\left(-\overrightarrow{BC}\times\overrightarrow{BC}\right)&=\mathbf 0\\ \left(\mathbf a\times\overrightarrow{BC}\right)+\left(-\overrightarrow{BC}\times\mathbf a\right)&=\mathbf 0\\ \left(\mathbf a\times\overrightarrow{BC}\right)+\left(\mathbf a\times\overrightarrow{BC}\right)&=\mathbf 0\\ 2\left(\mathbf a\times\overrightarrow{BC}\right)&=\mathbf 0\\ \mathbf a\times\overrightarrow{BC}&=\mathbf 0\\ \overrightarrow{OA}\times\overrightarrow{BC}&=\mathbf 0\end{align*}$\therefore$ $\overrightarrow{OA}$ and $\overrightarrow{BC}$ are parallel. (Shown)
\par\smallskip\noindent\textbf{(b)}\par\nopagebreak{}Since $\angle OBC=60^\circ$,\begin{align*}\overrightarrow{BO}\cdot\overrightarrow{BC}&=\left|\overrightarrow{BO}\right|\left|\overrightarrow{BC}\right|\cos60^\circ\\ \overrightarrow{BO}\cdot\overrightarrow{BC}&=\left|\overrightarrow{BO}\right|\left[2\left|\overrightarrow{OA}\right|\right]\cos60^\circ\\ -\mathbf b\cdot(\mathbf c-\mathbf b)&=(1)(2)(1)(0.5)\\ -\mathbf b\cdot\mathbf c+\mathbf b\cdot\mathbf b&=1\\ -\mathbf b\cdot\mathbf c+|\mathbf b|^2&=1\\ -\mathbf b\cdot\mathbf c+1^2&=1\\ \mathbf b\cdot\mathbf c&=0\end{align*}$\therefore$ $\overrightarrow{OB}$ and $\overrightarrow{OC}$ are perpendicular. (Shown)
\par\smallskip\noindent\textbf{(c)}\par\nopagebreak{}\begin{center}\begin{tikzpicture}[scale=0.6]\coordinate (O) at (0,0);\coordinate (H) at (0,3.2);\coordinate (B) at (2.6,0.6);\coordinate (C) at (2.9,-1.6);\draw (H)--(O)--(B)--(H)--(C)--(O);\draw (B)--(C);\node[left] at (O){\small$O$};\node[above] at (H){\small$H$};\node[right] at (B){\small$B$};\node[below] at (C){\small$C$};\end{tikzpicture}\end{center}$\mathbf h=\mathbf b\times\mathbf c\Rightarrow\mathbf h\perp\mathbf b$ and $\mathbf h\perp\mathbf c$\par $\left|\overrightarrow{OH}\right|$ is the height of the pyramid $HOBC$.\begin{align*}|\mathbf h|&=|\mathbf b\times\mathbf c|\\ &=|\mathbf b||\mathbf c|\sin90^\circ=(1)(\sqrt3)(1)=\sqrt3\quad(\angle BOC=90^\circ\ since\ \mathbf b\cdot\mathbf c=0)\\ \text{Area of }\Delta OBC&=\tfrac12(|\mathbf b\times\mathbf c|)=\tfrac12(|\mathbf h|)=\frac{\sqrt3}{2}\end{align*}Volume of pyramid $HOBC$\begin{align*}&=\left(\tfrac13\right)(\text{Area of }\Delta OBC)\left(\left|\overrightarrow{OH}\right|\right)=\left(\tfrac13\right)\left(\tfrac{\sqrt3}{2}\right)\left(\sqrt3\right)=\tfrac12\text{ units}^2\end{align*}
\par\medskip\hrule\medskip
\Q{\textcolor{black!60}{\small[EJC 2025 Prelim P1 Q6]}\par\smallskip There are three distinct points $A$, $B$ and $P$. The point $P$ lies on the circle with centre $O$ and diameter $AB$. Relative to the origin $O$, the position vectors of points $A$ and $P$ are $\mathbf a$ and $\mathbf p$ respectively.\begin{enumerate}[label=\textbf{(\alph*)},leftmargin=2.2em,itemsep=3pt]\item By expressing $\overrightarrow{AP}$ and $\overrightarrow{BP}$ in terms of $\mathbf a$ and $\mathbf p$, show that $\overrightarrow{AP}.\overrightarrow{BP}=0$.\hfill[3]\end{enumerate}The point $C$ with position vector $\mathbf c$ lies on $AB$ produced such that $AC:AB$ is $\lambda:1$.\begin{enumerate}[label=\textbf{(\alph*)},leftmargin=2.2em,itemsep=3pt,start=2]\item Find $\mathbf c$ in terms of $\lambda$ and $\mathbf a$.\hfill[2]\item It is given that $PC$ is a tangent to the circle and that angle $AOP$ is $120^\circ$. Using a suitable scalar product, find the value of $\lambda$.\hfill[4]\end{enumerate}}
\par\smallskip\noindent\textbf{Solution}\par \par\smallskip\noindent\textbf{(a)}\par\nopagebreak{}\begin{center}\begin{tikzpicture}[scale=0.6,>=Stealth]\draw (0,0) circle (1.6);\draw (-1.6,0)--(1.6,0);\draw[->,thick](0,0)--(-0.9,0) node[below]{\small$\mathbf a$};\draw[->,thick](0,0)--(1.13,1.13);\node at (0.45,0.75){\small$\mathbf p$};\node[left] at (-1.6,0){\small$A$};\node[right] at (1.6,0){\small$B$};\node[above right] at (1.13,1.13){\small$P$};\node[below] at (0,0){\small$O$};\end{tikzpicture}\end{center}\begin{align*}\overrightarrow{AP}&=\mathbf p-\mathbf a\\ \overrightarrow{BP}&=\mathbf p+\mathbf a\\ \overrightarrow{AP}.\overrightarrow{BP}&=(\mathbf p-\mathbf a).(\mathbf p+\mathbf a)\\ &=|\mathbf p|^2-|\mathbf a|^2\\ &=|\mathbf a|^2-|\mathbf a|^2\quad(\text{since $OA$ and $OP$ are radius, }\therefore|\mathbf p|=|\mathbf a|)\\ &=0\end{align*}
\par\smallskip\noindent\textbf{(b)}\par\nopagebreak{}\begin{center}\begin{tikzpicture}[scale=0.6,>=Stealth]\draw (0,0) circle (1.6);\draw[thick] (-1.6,0)--(2.6,0);\draw[->,thick](0,0)--(-0.9,0) node[above]{\small$\mathbf a$};\draw[->,thick](0,0)--(0.8,1.39) node[midway,left]{\small$\mathbf p$};\draw (0.8,1.39)--(2.6,0);\node[left] at (-1.6,0){\small$A$};\node[below right] at (1.6,0){\small$B$};\node[below right] at (2.6,0){\small$C$};\node[above] at (0.8,1.39){\small$P$};\node[below] at (0,0){\small$O$};\node at (-0.8,-0.5){\small$1$};\node at (2.1,-0.5){\small$\lambda-1$};\end{tikzpicture}\end{center}By Ratio Theorem,\begin{align*}\mathbf b&=\frac{(\lambda-1)\mathbf a+\mathbf c}{\lambda}\\ \mathbf c&=\lambda\mathbf b-(\lambda-1)\mathbf a\\ \mathbf c&=\lambda(-\mathbf a)-(\lambda-1)\mathbf a\\ &=(1-2\lambda)\mathbf a\end{align*}
\par\smallskip\noindent\textbf{(c)}\par\nopagebreak{}\begin{align*}\overrightarrow{PC}&=(1-2\lambda)\mathbf a-\mathbf p\\ \overrightarrow{OP}.\overrightarrow{PC}&=0\\ \mathbf p.((1-2\lambda)\mathbf a-\mathbf p)&=0\\ (1-2\lambda)\mathbf p.\mathbf a-\mathbf p.\mathbf p&=0\\ (1-2\lambda)|\mathbf p||\mathbf a|\cos120^\circ-|\mathbf p|^2&=0\\ -\tfrac12(1-2\lambda)|\mathbf a|^2-|\mathbf a|^2&=0\\ -\tfrac12(1-2\lambda)&=1\\ \lambda&=\frac32\end{align*}
\par\medskip\hrule\medskip
\par\vspace{\baselineskip}\begin{center}\textbf{END OF WORKSHEET}\end{center}\end{document}
